\documentclass[12pt]{article}

% --- Core math + proof tools ---
\usepackage{amsmath, amssymb, amsthm}

% --- Lists and spacing ---
\usepackage{enumitem}
\setlist[enumerate,1]{label=\textbf{(\alph*)}, leftmargin=2em, itemsep=0.25em}
\setlist[itemize]{leftmargin=1.5em, itemsep=0.25em}
\setlength{\parskip}{0.35em}
\setlength{\parindent}{0pt}

% --- Page geometry + header/footer ---
\usepackage{geometry}
\geometry{margin=1in}
\usepackage{csquotes} % to use blockquotes

\usepackage{fancyhdr}
\pagestyle{fancy}
\fancyhf{}
\rhead{MATH 421 — 003, Fall 2025}
\lhead{Homework 6}
\cfoot{\thepage}
\renewcommand{\headrulewidth}{0.4pt}

% --- Links (nice for references, stays subtle) ---
\usepackage[hidelinks]{hyperref}

% --- Title info ---
\title{Homework 6}
\author{Alejandro De La Torre \\ MATH 421 — The Theory of Single Variable Calculus \\ Section 003 \\ Instructor: Grace Work}
\date{October 2025}

% --- (Optional) proof environments you might want ---
\newtheorem*{claim}{Claim}
\newtheorem*{lemma}{Lemma}
\newtheorem*{remark}{Remark}

% --- Convenience: a Problem header command ---
\newcommand{\problem}[2][]{%
  \vspace{0.5em}
  \hrule\vspace{0.4em}
  \textbf{Problem #2}\if\relax\detokenize{#1}\relax\else\; \textit{(#1)}\fi
  \vspace{0.25em}\hrule\vspace{0.8em}
}


\begin{document}
\maketitle

% ------------------ collaboration + sources ------------------
\section*{Collaboration Statement}
%TODO


\bigskip

% ------------------ problems ------------------

\problem[ Ross Chapter 2, 9.3 ]{1}
\textit{Suppose $\lim a_n = a$, $\lim b_n = b$, and $s_n = \dfrac{a_n^3 + 4a_n}{b_n^2 + 1}$. Prove that $\lim s_n = \dfrac{a^3 + 4a}{b^2 + 1}$ carefully, using the limit theorems.}

\vspace{0.75em}
\textbf{Discussion.} \\%(Very brief plan / key ideas.)
Here we essentially want to break the problem up into the individual limit components and apply the limit theorems 9.1-9.7 accordingly to show that the limit converges to $\frac{a^3+4a}{b^2+1}$.

For reference, here is what each of the 9.1-9.7 theorems say in a nutshell:
\begin{center}
\begin{minipage}{0.7\textwidth}
\begin{footnotesize}
\begin{enumerate}[label=9.\arabic* Theorem:]
    \item Convergent sequences are bounded.
    \item If $\lim s_n = s$, then $\lim (ks_n) = k \lim s_n$.
    \item If $\lim s_n = s$ and $\lim t_n = t$, then $\lim (s_n + t_n ) = \lim s_n + \lim t_n$.
    \item If $\lim s_n = s$ and $\lim t_n = t$, then $\lim (s_n t_n) = (\lim s_n)(\lim t_n)$.
    \item If $\lim s_n = s$, $(s \neq 0)$  and $s_n \neq 0$ for all n, then $\lim \frac{1}{s_n}= \frac{1}{s}$.
    \item If $\lim s_n = s$, $(s \neq 0)$, $s_n \neq 0$ for all n, and $\lim t_n = t$, then $\lim \frac{t_n}{s_n}=\frac{t}{s}$.
    \item  
    \begin{enumerate}
        \item $\lim \frac{1}{n^p}=0$ for $p>0$.
        \item $\lim a^n=0$ if $|a|<1$.
        \item $\lim (n^{\frac{1}{n}})=1$.
        \item $\lim (a^{\frac{1}{n}})=1$ for $a > 0$.
    \end{enumerate}
\end{enumerate}
\end{footnotesize}
\end{minipage}
\end{center}
\vspace{0.25em}

\textbf{Proof.}\\
Given that $\lim a_n = a$ and $\lim b_n = b$, we can then express $\lim\frac{a_n^3 + 4a_n}{b_n^2 + 1}$ as $\frac{\lim a_n^3 + \lim 4a_n}{\lim b_n^2 + \lim 1}$. We can express the result of $\lim a_n^3$ as $a^3$, and  $\lim b_n^2=b^2$ by theorem 9.4. Next we can show that $\lim 4a_n = 4a$ by theorem 9.2. Then, using the result from theorem 9.3, we can show that $\lim (a_n^3+4a_n) = a^3+4a$ and $\lim (b_n^2+1) = b^2+1$. Finally, using the result from theorem 9.6, we show that:
\[
\lim s_n = \frac{\lim (a_n^3 + 4a_n)}{\lim (b_n^2 + 1)} = \frac{a^3+4a}{b^2+1}
\]
for $b^2+1 \neq 0$ for all n.
\qed

\problem[Ross Chapter 2, 9.9(b)]{2}
% The prompt asks you to convince yourself that (a) and (c) are true (no write-up needed).

\textit{Suppose there exists $N_0$ such that $s_n \leq t_n$ for all $n > N_0$. 
Prove that if $\lim t_n = -\infty$, then $\lim s_n = -\infty$.}

\vspace{0.75em}

\textbf{Discussion.} \\
Here we want to leverage the definition of divergence which is given in Ross 9.8 theorem, which states that:

\begin{quote}
    \textit{We write $\lim s_n = -\infty$ provided for each $M < 0$ there is a number such that $ n > N$ implies $s_n < M$.}
\end{quote}

The nuance here is that we are comparing a sequence diverging to a 'negative infinity' ($t_n$) to another sequence that we are given to suppose is 'more negative' ($s_n$) after a certain point ($N_0$;  for all $n > N_0$). In other words, we are comparing a sequence diverging to negative infinity ($t_n$) with another sequence ($s_n$) that is eventually less than or equal to it.

So if $\lim t_n = -\infty$ is true, then we know by theorem 9.8 that for every $M<0$ there is a number $N_1$ such that $n > N_1$ implies $t_n < M$. We also know that there has got to be a $N_0$ (we suppose there is) such that $s_n \leq t_n$ for all $n > N_0$, so we can then just define an $N$ that is greatest of the set of $\{N_0, N_1\}$ in whatever arbitrary case we take on (which we know much exist) so that we can say $s_n \leq t_n \leq M$ so that we can henceforth conclude that $s_n < M$ which is what we need to say that $s_n$ diverges.
\vspace{0.25em}

\textbf{Proof.}\\
Let $M<0$. We know that there exists an $N_1$ such that $n > N_1$ implies $t_n \leq M$.

Let $N = \max \{N_0, N_1\}$. Thus, $n > N$ implies
\[
s_n \leq t_n \leq M
\]
which implies
\[
s_n \leq M
\]
Therefore $\lim s_n = - \infty$
\qed

\problem[ Ross Chapter 2, 9.10(b) ]{3}
% Again, only part (b) is to be written; (a) and (c) you just check mentally.

\textit{Show that $\lim s_n = +\infty$ if and only if $\lim (-s_n) = -\infty$.}
\vspace{0.75em}

\textbf{Discussion.}
We want to utilize theorem 9.8 and use a distinct M according to each case (since this is an iff statement) and prove divergence accordingly.

We will begin by show it first in the forward ($\implies$) direction: "if $\lim s_n = + \infty$, then $\lim (-s_n) = -\infty$, wherein we would let $M < 0$ so that we can say $ s_n > -M$ then manipulate the expression to arrive at the conclusion that if $\lim s_n = +\infty$ then $\lim (-s_n) = - \infty$.

Then we would show the backward ($\impliedby$) direction of the iff statement such that: "if $\lim (-s_n) = -\infty$ then $\lim s_n = + \infty$", and follow the same procedure but with $M>0$ to show the divergence toward a positive $\infty$ result.
\vspace{0.25em}

\textbf{Proof.}\\
    ($\implies$) Let $\lim s_n = - \infty$. Let $M < 0$. For each $M < 0$ there is a number $N$ such that $n > N$ implies
    \[
    s_n > -M
    \]
    which implies
    \[
    -s_n < M
    \]
    Thus, $\lim (-s_n) = -\infty$ \\ 

\bigskip

    ($\impliedby$) Let $\lim(-s_n) = - \infty$. Let $M > 0$, For each $M > 0$ there us a number $N$ such that $n>N$ implies 
    \[
    (-s_n) > -M
    \]
    which implies
    \[
    s_n > M
    \]
    Thus, $\lim(s_n) = \infty$
\qed

\problem[ Ross Chapter 2, 10.2 ]{4}

\textit{Prove Theorem 10.2 for bounded decreasing sequences.}
\vspace{0.75em}

\textbf{Discussion.}\\
Theorem 10.2 states that \textit{“All bounded monotone sequences converge.”} In the textbook, this theorem is proven first for the case of a bounded increasing sequence, using the concept of the supremum (\textit{sup}). To prove the theorem for a bounded decreasing sequence, we instead use the \textit{infimum} (\textit{inf}) as our central idea. Intuitively, the argument mirrors the increasing case: rather than the terms of the sequence being squeezed upward toward the supremum, here they are squeezed downward toward the infimum.

So, given a bounded decreasing sequence $(s_n)$, we recognize that since it is bounded below, there exists an $\inf S$, say $L$, which acts as the “greatest lower bound” of the set $\{s_n : n \in \mathbb{N}\}$. The proof strategy is to show that $s_n$ converges to this infimum $L$ using the $\varepsilon$-definition of convergence, just as the increasing case shows convergence to the supremum. The main difference is in the direction of the inequalities: where the increasing case used $s_N > u - \varepsilon$, we will use $s_N < L + \varepsilon$. This reflects how a decreasing sequence approaches its lower bound from above.

\vspace{0.25em}

\textbf{Proof.}\\
Let $(s_n)$ be a bounded decreasing sequence. Since $(s_n)$ is bounded, the set $S = \{s_n : n \in \mathbb{N}\}$ is bounded below. By the Completeness Axiom, there exists an infimum $L = \inf S$.

Let $\varepsilon > 0$. Since $L + \varepsilon$ is not a lower bound for $S$, there exists some $N \in \mathbb{N}$ such that $s_N < L + \varepsilon$. Because the sequence is decreasing, for all $n > N$ we have
\[
s_n \leq s_N < L + \varepsilon.
\]
We also know that, by definition of the infimum, $L \leq s_n$ for all $n$. Combining these inequalities, we get
\[
L \leq s_n < L + \varepsilon,
\]
which implies that
\[
|s_n - L| < \varepsilon.
\]
Therefore, $\lim s_n = L = \inf S$.

\qed

\problem[ Ross Chapter 2, 10.5 ]{5}

\textit{Prove Theorem 10.4(ii).}

\textbf{Discussion.}\\
Theorem 10.4 (ii) states:

\begin{quote}
\textit{If $(s_n)$ is an unbounded decreasing sequence, then $\lim s_n = -\infty$.}
\end{quote}

To prove this, we should recall the definitions of an \textit{unbounded} and a \textit{decreasing} sequence.

A sequence $(s_n)$ is called \textbf{decreasing} if for every $n \in \mathbb{N}$, $s_{n+1} \leq s_n$.  
It is called \textbf{unbounded below} if for every real number $M < 0$, there exists an $N \in \mathbb{N}$ such that $s_N < M$.  

Together, these properties describe a sequence whose terms continually decrease without bound. Our goal is to show that this precisely means the limit of $(s_n)$ diverges to negative infinity.

The strategy of the proof parallels that of part (i) of Theorem 10.4, which treats the unbounded increasing case (where $\lim s_n = +\infty$). Instead of showing the sequence eventually exceeds any given positive number, we will show that for every negative number $M$, the sequence eventually falls below it. We will also use the monotonic property (that $s_n$ is decreasing) to ensure that once the sequence passes below $M$, it stays there for all subsequent $n$.

\vspace{0.25em}

\textbf{Proof.}\\
To show that if $(s_n)$ is an unbounded decreasing sequence, then $\lim s_n = -\infty$, let $M < 0$.

Since the set $\{s_n : n \in \mathbb{N}\}$ is unbounded and it is bounded above by $s_1$, it must be unbounded below. Hence, there exists some $N \in \mathbb{N}$ such that $s_N < M$.

Because $(s_n)$ is decreasing, for all $n > N$ we have
\[
s_n \leq s_N < M.
\]
Thus, for every $M < 0$, there exists an $N$ such that $n > N$ implies $s_n < M$.

By definition of divergence to negative infinity, we conclude that
\[
\lim s_n = -\infty.
\]

\qed

\problem[ Ross Chapter 2, 10.7 ]{6}
\textit{Let $S$ be a bounded nonempty subset of $\mathbb{R}$ such that $\sup S$ is not in $S$. Prove there is a sequence $(s_n)$ of points in $S$ such that $\lim s_n = \sup S$. See also Exercise 11.11.}

\vspace{0.75em}

\textbf{Discussion.} \\
I am feeling really slow right now, so I am going to take things very slowly for this one. So let's begin by restate what we know. We know:

    \begin{enumerate}[label=(\arabic*)]
        \item $S$ is bounded. (So $\sup S$ exists by the Completeness Axiom.)
            \begin{enumerate}
                \item $S$ being bounded means there exists some $M \in \mathbb{R}$ such that for all $s \in S$, $|s| \leq M$. So $[-M, M]$ is an interval that contains $S$.
                \item From the completeness axiom, we know every nonempty set of real numbers that is bounded above has a supremum in $\mathbb{R}$. In other words we know S has a supremum by the completeness axiom since it is bounded.
            \end{enumerate}
        \item $S$ is nonempty: $S \neq \emptyset$. (So there is at least one point in our sequence).
        \item $S \subseteq \mathbb{R}$.
        \item $\sup S \notin S$.
    \end{enumerate}

From these three points, we can then construct a sequence $(s_n)$ in $S$ that converges to $\sup S$. The thing is, we do not know what $\sup S$ is exactly (it's arbitrary), but we do know that it is the least upper bound of $S$. 

So to keep things arbitrary, we can define a $t$ such that $t = \sup S$ and since $S \subseteq \mathbb{R}$, we know $t \in \mathbb{R}$. We know $t$ is the least upper bound of $S$ and must indeed exist for any arbitrary set $S$ by the completeness axiom, but not an element of $S$.

Now, to construct our sequence $(s_n)$ in $S$ that converges to $t$, we can use the definition of supremum. Since $t$ is the least upper bound of $S$, for every $\varepsilon > 0$, there exists an element $s \in S$ such that $t - \varepsilon < s \leq t$. This means we can find elements of $S$ that are arbitrarily close to $t$, without $t$ itself being in $S$.

So we can define our sequence $(s_n)$ by choosing $s_n \in S$ such that $t - \frac{1}{n} < s_n \leq t$ for each natural number $n$. This way, as $n$ increases, $s_n$ gets closer and closer to $t$.

So that set up alone intuitively takes us to the conclusion that $\lim s_n = t = \sup S$. So this is awesome because now we have a sequence defined in a way that converges to $\sup S$. Half the battle is won. \\

The other half of the battle involves us proving convergence and then formalize this into a proof. Given our set up, we can easily prove convergence via the squeeze lemma (which we proved in \textit{Chapter 2 \S8, Exercise $8.5 (a)$}) which states that if we have three sequences $(a_n)$, $(b_n)$, and $(c_n)$ such that $a_n \leq b_n \leq c_n$ for all $n$ beyond some index $N$, and if $\lim a_n = \lim c_n = L$, then $\lim b_n = L$ as well. So check this proof out below.

\small Shoutout Professor Work for helping me with this one. My bad that my brain was exceptionally smooth when I came into office hours. \normalsize
\vspace{0.25em}

\textbf{Proof.} \\ 
Let $S \subseteq \mathbb{R}$ be a bounded nonempty set such that $\sup S \notin S$. Let $t = \sup S$. Since $t$ is the least upper bound of $S$, for every $\varepsilon > 0$, there exists an element $s \in S$ such that $t - \varepsilon < s \leq t$. \\

Define the sequence $(s_n)$ by choosing $s_n \in S$ such that $t - \frac{1}{n} < s_n \leq t$ $\forall n \in \mathbb{N}$. In other words, for each natural number $n$, we select an element $s_n$ from $S$ that lies within the interval $(t - \frac{1}{n}, t]$. \\

From the squeeze lemma, we can then say that since $t - \frac{1}{n} < s_n \leq t$, we can define two sequences $(a_n)$ and $(c_n)$ such that $a_n = t - \frac{1}{n}$ and $c_n = t$. We can then see that $a_n \leq s_n \leq c_n$ for all $n$. 
Now we can compute the limits of $(a_n)$ and $(c_n)$:
\[
\lim a_n = \lim (t - \frac{1}{n}) = t - 0 = t
\]
and
\[\lim c_n = \lim t = t
\]
By the squeeze lemma, since $\lim a_n = \lim c_n = t$, we have $\lim s_n = t = \sup S$.
\qed

\end{document}